\def\centerarc[#1](#2)(#3:#4:#5){ \draw[#1] ($(#2)+({#5*cos(#3)},{#5*sin(#3)})$) arc (#3:#4:#5); }
%% https://tex.stackexchange.com/questions/55068/is-there-a-tikz-equivalent-to-the-pstricks-ncbar-command
\tikzset{
    ncbar angle/.initial=90,
    ncbar/.style={
        to path=(\tikztostart)
        -- ($(\tikztostart)!#1!\pgfkeysvalueof{/tikz/ncbar angle}:(\tikztotarget)$)
        -- ($(\tikztotarget)!($(\tikztostart)!#1!\pgfkeysvalueof{/tikz/ncbar angle}:(\tikztotarget)$)!\pgfkeysvalueof{/tikz/ncbar angle}:(\tikztostart)$)
        -- (\tikztotarget)
    },
    ncbar/.default=0.5cm,
}

\tikzset{round left paren/.style={ncbar=0.5cm,out=120,in=-120}}
\tikzset{round right paren/.style={ncbar=0.5cm,out=60,in=-60}}

\newcommand{\quartercirc}{%
\begin{tikzpicture}
    \centerarc[black,fill=gray!30](0,0)(0:90:0.24);
    \filldraw[fill=gray!30,opacity=0.0,fill opacity=1] (0,0) -- (0,0.24) -- (0.24,0) -- cycle;
    \draw[-] (0,0) to (0,0.24);
    \draw[-] (0,0) to (0.24,0);
\end{tikzpicture}%
}


\begin{figure}[h!]
    \centering
    \begin{tikzpicture}[tips=proper]
        % HOW MUCH TO RIGHT INCREMENT LARGE DENSITY SQUARES
        \def\inc{1.5}
        % HOW MUCH TO RIGHT INCREMENT SMALL DENSITY SQUARES
        \def\sinc{1.2}
        % SIZE OF SMALL DENSITY SQUARES
        \def\s{0.75}
    
        \node[isosceles triangle,
            draw,
            thick,
            isosceles triangle apex angle=60,
            rotate=90,
            minimum size=6cm] (T1) at (0,0){};
        
        \draw[black,dashed] (T1.east) .. controls ([xshift=-0.1cm]T1.358) .. ([yshift=-0.5cm]T1.east) node[] (N1) {{}};
        \draw[black,dashed] ([yshift=-0.5cm]T1.east) .. controls ([xshift=0.4cm]T1.40) and (T1.350) .. (T1.center) node[fill=white,pos=0.9] (N1) {$\pi$};
        \draw[black,dashed] (T1.center) .. controls ([xshift=2.4cm]T1.110) .. (T1.west) node[] (N2) {{}};

        % LOCATION OF TOP HARDCORE DENSITIES
        \coordinate (top) at (5,4);
        \coordinate (top2) at ($(top)+(\inc,0)$); 
        \coordinate (top3) at ($(top2)+(\inc,0)$);
        
        \node[rectangle,
    	draw = black,
    	minimum width = 1cm, 
    	minimum height = 1cm] (r) at (top) {};
        \node[ellipse,
            draw = black,
            fill=gray!30,
            minimum width = 0.4cm,
            minimum height = 0.4cm] (e) at (top) {\footnotesize $H$};
        \node[rectangle,
    	draw = black,
    	minimum width = 1cm, 
    	minimum height = 1cm] (r2) at (top2) {};
        \node[ellipse,
            draw = black,
            fill=gray!30,
            minimum width = 0.5cm,
            minimum height = 0.5cm] (e2) at (top2) {\footnotesize $H$};
        \node[rectangle,
    	draw = black,
    	minimum width = 1cm, 
    	minimum height = 1cm] (r3) at (top3) {};
        \node[ellipse,
            draw = black,
            fill=gray!30,
            minimum width = 0.5cm,
            minimum height = 0.5cm] (e3) at (top3) {\footnotesize $H$};

        % TOP PARENTHESES
        \draw [black] (4.3,3.4) to [round left paren ] (4.3,4.6);
        \draw [black] (8.7,3.4) to [round right paren ] (8.7,4.6);

        % TOP COMMAS
        \draw [black] (5.75,3.5) node [] {\LARGE ,};
        \draw [black] (7.25,3.5) node [] {\LARGE ,};

        % BOTTOM PARENTHESES
        \draw [black] (4.8,-2.5) to [round left paren ] (4.8,-1.5);
        \draw [black] (8.4,-2.5) to [round right paren ] (8.4,-1.5);

        % BOTTOM COMMAS
        \draw [black] (6,-2.4) node [] {\LARGE ,};
        \draw [black] (7.2,-2.4) node [] {\LARGE ,};

        % LOCATION OF BOTTOM HARDCORE DENSITIES
        \coordinate (bot) at (5.4,-2); 
        \coordinate (bot2) at ($(bot)+(\sinc,0)$);
        \coordinate (bot3) at ($(bot2)+(\sinc,0)$);
        
        \node[rectangle,
    	draw = black,
    	minimum width = \s cm, 
    	minimum height = \s cm] (r2) at (bot) {};
        \node[rectangle,
    	draw = black,
    	minimum width = \s cm, 
    	minimum height = \s cm] (r3) at (bot2) {};
        \node[rectangle,
    	draw = black,
    	minimum width = \s cm, 
    	minimum height = \s cm] (r4) at (bot3) {};
       
        \centerarc[black,fill=gray!30]([xshift=-0.38cm,yshift=-0.38cm]bot)(0:90:0.5);
        \centerarc[black,fill=gray!30]([xshift=0.38cm,yshift=0.38cm]bot2)(180:270:0.5);
        \centerarc[black,fill=gray!30]([xshift=0.38cm,yshift=-0.38cm]bot3)(90:180:0.5);

        % FILLING IN LIGHT GRAY REST OF BLOCKS
        \filldraw[fill=gray!30,opacity=0.0,fill opacity=1] (r2.225) -- (r2.162) -- (r2.288) -- cycle;
        \filldraw[fill=gray!30,opacity=0.0,fill opacity=1] (r3.342) -- (r3.45) -- (r3.108) -- cycle;
        \filldraw[fill=gray!30,opacity=0.0,fill opacity=1] (r4.251) -- (r4.315) -- (r4.18) -- cycle;

        % path label
        \draw[color=black] (T1.west) node [below,fill=white] {\footnotesize Leaf $\ell$};
        \node[draw,circle,fill=black,inner sep=1pt] (x) at (T1.west) {};

        % LABELS
        % \draw[-{Stealth[scale=1]}] (6,2.7) to (e);
        % \draw [black] (6,2.7) node [fill=white] {Hardcore measure $H$};

        \draw[-{Stealth[scale=1]}] (4.6,-3.3) to ([xshift=-.2 cm,yshift=-.2cm]bot);
        \draw [black] (4.6,-3.3) node [fill=white] {Hardcore measure $H_{\pi_1}$ with $\mathrm{Dens}_H(\ell, 1)={\mathrm{area}(\quartercirc)}/{\mathrm{area}(\square})$};

        % QUERIES ALONG PATH PI
        \draw[-{Stealth[scale=1]}] (6.5,3) to node[midway,fill=white!30,scale=1] {Queries along path $\pi$} (6.5,-1);

       % REDRAWS
        \node[rectangle,
    	draw = black,
    	minimum width = \s cm, 
    	minimum height = \s cm] (aaaa) at (bot) {};
        \node[rectangle,
    	draw = black,
    	minimum width = \s cm, 
    	minimum height = \s cm] (bbbbb) at (bot2) {};
        \node[rectangle,
    	draw = black,
    	minimum width = \s cm, 
    	minimum height = \s cm] (ccccc) at (bot3) {};
        
    \end{tikzpicture}
    \captionsetup{width=.9\linewidth}
    \caption{Illustration of a hardcore density. The tree $T:(\bn)^3\to\bits^3$ seeks to compute a function $f^{\otimes 3}$. The tuple of squares at the top of the figure illustrates the set of all inputs to the function while the strings in the support of the hardcore measure are shaded gray. The tuple at the bottom of the figure illustrates the set of inputs reaching the leaf $\ell$. Each block is the subcube consistent with the path $\pi$ and the shaded region denotes the fragment of $H$ which is contained in the corresponding subcube. }
    \label{fig:hardcore}
\end{figure}
